\begin{abcliste}
\abc
Falsch ist \result{E_\mathrm{el} = \EaP{0}{3} \text{ in (3)}}. In (3) ist die Feder entspannt, die Spannenergie also null: $E_\mathrm{el} = \frac{1}{2}\,k\cdot 0^2 = 0$. Die ganze Spannenergie ist zu kinetischer Energie geworden. So, wie der Schüler es eingetragen hat, wäre die Gesamtenergie in (3) doppelt so gross wie am Anfang.

\abc
Die Spannenergie wächst mit dem Quadrat der Zusammendrückung:
\begin{align}
 E_\mathrm{el,2} &= \EebF = \frac{1}{4}\cdot\frac{1}{2}\,k\,s^2 = \Eeb \approx \EebP{0}{3}
\end{align}
Die Feder gibt also schon in der ersten Hälfte des Wegs drei Viertel ihrer Energie ab, weil sie dort stärker drückt:
\begin{align}
 E_\mathrm{kin,2} &= \EkbF = \Ekb \approx \EkbP{0}{3}
\end{align}
Die Zeile (2) stimmt: Die Summe ist wie in (1) \EaP{0}{3}.

\abc
In (3) ist $\frac{1}{2}\,m\,v^2 = \frac{1}{2}\,k\,s^2$:
\begin{align}
 v_3 &= \vcF = \sK\cdot\sqrt{\frac{\kK}{\mK}} = \vc \approx \result{\vcP{0}{2}}
\end{align}
In (2):
\begin{align}
 v_2 &= \vbF = \sqrt{\frac{2\cdot\Ekb}{\mK}} = \vb \approx \result{\vbP{0}{2}}
\end{align}
Nach der halben Strecke hat der Klotz schon $\sqrt{3}/2 \approx \SI{87}{\percent}$ seiner Endgeschwindigkeit.

\abc
$\frac{1}{2}\,k\,s = \FwP{0}{3}$ hat die Einheit $\si{\newton\per\meter}\cdot\si{\meter} = \si{\newton}$: Das ist eine Kraft (die mittlere Federkraft beim Entspannen), keine Energie. Die Spannenergie ist $\frac{1}{2}\,k\,s^2$, mit der Einheit $\si{\newton\meter} = \si{\joule}$.
\end{abcliste}
